\documentclass{article}
\usepackage{pathfinder-paper}

\title{Where Comparison Rankings Enter Open-Team Reward Shaping}

\begin{document}
\maketitle

\begin{abstract}
Q asks that a prediction guarantee be connected to a downstream objective and that predictor calls be
charged; P turns multimodal pairwise comparisons into agent-specific potential-based rewards.
% ledger: 1, 2, 12, 19
We analyse the interface between these claims, without asserting a new return guarantee. With complete
per-agent reward accounting, P's shaping telescopes from exact discounted returns and cancels from a matched
critic's temporal-difference residual. With full-population pooling, its softmax normalisation also removes
every individual ranking from the pooled shaping increment. An active-only reward mask instead leaves a
reentry term and can preserve ranking dependence at coalition changes. These are conditional consequences of
the stated potential, not findings about P's implementation. % ledger: 4, 5, 9, 10, 14, 15, 16, 17, 19
A two-agent sample additionally shows that connectedness alone does not ensure the finite unregularised
Bradley--Terry estimate required by P's stated error bound. The remaining question is whether comparisons
improve finite-budget unshaped performance through a measured training channel after their calls are charged.
% ledger: 3, 5, 8, 12, 19
\end{abstract}

\section{The two papers and the question}

Q surveys learning-augmented algorithms and distinguishes predictor error from the downstream cost it is
meant to control. Its inference-graph account requires a justified edge between components, and its
prediction-pricing proposal explicitly counts calls in objective units~\cite{Q}. % ledger: 1, 2, 4, 12, 19
P's MARS-RA queries a multimodal comparator for each ordered pair in the active coalition, fits
Bradley--Terry scores, softmaxes them into a vector potential, and adds the potential difference to a scalar
environment reward for multi-agent training~\cite{P}. It reports task-success and comparison-accuracy trends;
its estimation proposition addresses latent comparator scores, while its Shapley proposition assumes
alignment of those scores with Shapley values. Neither proposition states a bound on the return of the
learned policy per comparison call. % ledger: 3, 8, 12, 19

We ask which reward or critic route could carry P's ranking signal into a finite-training improvement when
Q's downstream and cost requirements are applied. The calculations below identify routes to inspect. They do
not establish that any route is used by P's implementation. % ledger: 12, 15, 17, 19

\section{Conditional accounting for P's potential}

Let \(\mathcal N\) be the fixed agent population, \(I^t\subseteq\mathcal N\) the active coalition, and
\(a_{i,t}=\mathbf 1\{i\in I^t\}\). For a nonempty nonterminal coalition, let \(\psi_{i,t}\) be P's softmax
score for \(i\in I^t\), embedded as zero for an inactive agent. Use one cached value at each endpoint of
adjacent transitions and set \(\psi_{i,T}=0\) at termination. Write
\(F_{i,t}=\gamma\psi_{i,t+1}-\psi_{i,t}\), with shaping weight \(\rho\). This fixed-coordinate,
cached-potential convention is a condition of our analysis; P does not specify its cross-coalition reward
convention. % ledger: 4, 5, 6, 10, 15, 16

\begin{proposition}[Complete accounting and active masking]
If agent \(i\)'s shaping increment is credited on every transition, then, pathwise,
\[
\sum_{t=0}^{T-1}\gamma^t F_{i,t}=-\psi_{i,0}.
\]
If it is credited only when \(a_{i,t}=1\), then
\[
\sum_{t=0}^{T-1}\gamma^t a_{i,t}F_{i,t}
=-a_{i,0}\psi_{i,0}
+\sum_{t=1}^{T-1}\gamma^t(a_{i,t-1}-a_{i,t})\psi_{i,t}.
\]
At a reentry time \(t\), the latter sum includes \(-\gamma^t\psi_{i,t}\).
\end{proposition}
% ledger: 4, 5, 6, 10, 16, 19
\begin{proof}
In the first sum, every intermediate \(\psi_{i,t}\) appears once with each sign, leaving
\(-\psi_{i,0}+\gamma^T\psi_{i,T}\). In the masked sum, the coefficient of \(\psi_{i,t}\), for \(1\leq t<T\),
is \(\gamma^t(a_{i,t-1}-a_{i,t})\); the terminal term is zero. % ledger: 4, 10, 16
\end{proof}

The first identity is the standard potential-shaping cancellation applied to P's coordinates. For a fixed
initial potential determined before the first action, its pathwise offset does not change exact policy
ordering. The second identity shows why zeroing inactive potentials is insufficient when rewards are masked:
the transition into a returning agent is not credited to that agent. Neither identity says how P actually
routes its rewards. % ledger: 5, 6, 10, 12, 16

\begin{proposition}[Pooling fork]
Let \(u_t=\sum_{i\in\mathcal N}\psi_{i,t}\). If every coordinate is pooled, the shaping increment is
\(\sum_{i\in\mathcal N}F_{i,t}=\gamma u_{t+1}-u_t\). P's softmax gives \(u_t=1\) for each nonempty
nonterminal coalition, so this pooled increment is independent of the individual rankings. If instead only
\(i\in I^t\) are pooled and both endpoint coalitions are nonempty with a nonterminal successor, then
\[
\sum_{i\in I^t}F_{i,t}
=\gamma\left(1-\sum_{i\in I^{t+1}\setminus I^t}\psi_{i,t+1}\right)-1.
\]
Thus entrant scores can affect this active-only pool.
\end{proposition}
% ledger: 6, 14, 15, 16, 17
\begin{proof}
Sum the potential differences, then use \(\sum_{i\in I^t}\psi_{i,t}=1\). In the active-only case, the
successor sum includes exactly the agents active at both endpoints, leaving out entrants. % ledger: 14, 15, 17
\end{proof}

The proposition concerns the shaping increment, rather than a claim that P uses a pooled team reward.
Separate agent rewards can still change approximate training updates, even where complete per-agent returns
telescope. % ledger: 14, 15, 17, 19

For a critic \(\widehat V_i^{\mathrm{env}}\), set \(\widetilde V_i=\widehat V_i^{\mathrm{env}}-\rho\psi_i\).
Under complete accounting its shaped temporal-difference residual is exactly the unshaped residual. If the
implemented critic adds mismatch \(e_i\), the residual difference is \(\gamma e_{i,t+1}-e_{i,t}\). The same
cancellation holds for a generalised-advantage sum formed from matched residuals. This identifies critic
mismatch as a possible training channel, not a bound on final return. % ledger: 5, 9, 12, 19

\section{Two limits on transferring an accuracy claim}

\begin{proposition}[Connected comparisons need not yield a finite estimate]
For two agents at a true Bradley--Terry tie, take two independent comparisons of their sole pair. The
undirected comparison graph is connected, but the unregularised maximum-likelihood estimate is not finite on
the event that both outcomes favour the same agent. This event has probability \(1/2\).
\end{proposition}
% ledger: 3, 5, 12
\begin{proof}
Each ordered draw favours either agent with probability \(1/2\). Each unanimous outcome has probability
\(1/4\). For either outcome, increasing the winning agent's score difference strictly increases the
likelihood, whose supremum is reached only at an infinite difference. % ledger: 3, 5
\end{proof}

P states a finite score-error guarantee with probability at least \(1-n^{-2}=3/4\) under connectedness. The
example permits a finite estimate on at most half the samples, so connectedness by itself cannot support that
assertion for its unregularised fit. This is an objection to that stated estimator premise, not to P's
reported experiments; it does not assess a constrained or regularised estimator. % ledger: 3, 5, 12, 19

A distinct one-step diagnostic also limits what pairwise accuracy can imply. If agent 2 is inert, agent 1
chooses \(a\sim\operatorname{Bernoulli}(1/2)\), \(G=a\), and the logit score is \(z=a-1/2\), then the
baseline estimator \(g_b=z(G-b)\) has mean \(1/4\) and variance \((b-1/2)^2/4\). The neutral \(b=1/2\) has
zero variance; a rank-correct softmax of Shapley scores \((1/2,0)\) gives \(b>1/2\) and positive variance.
This example tests a proposed downstream error measure, not MAPPO performance. % ledger: 7, 11, 12

\section{Prior work and interpretation}

The central cancellation is published potential-shaping theory, rather than a new invariance theorem. P
itself cites dynamic shaping work; Gupta and colleagues show that potential shaping can leave expected
policy-gradient updates unchanged while increasing their variance~\cite{gupta2023}. Müller and Kudenko
already obtain a shaping potential from vision-language pairwise preferences and study label accuracy and
learning speed in single-agent environments~\cite{muller2026}. Thus using visual preferences for potential
shaping, and warning that ranking quality alone does not establish gradient quality, are existing steps.
% ledger: 6, 9, 11, 12, 19; related-work verification: arXiv:2606.27180

What this note contributes for the Q--P pair is the explicit three-way accounting fork: complete per-agent
rewards telescope, full-population pooling erases P's softmax ranking, and active-only credit can expose
entrant scores while leaving a reentry term. The fork is algebraic and conditional. It does not establish
novelty over all reward-routing implementations, nor does it identify P's implementation.
% ledger: 10, 14, 15, 16, 17, 19

P's stated all-ordered-pairs protocol calls the comparator \(q_t=|I^t|(|I^t|-1)\) times at a queried step. P
reports 50-million-step training times of 16 hours for MAPPO and 30 hours for MARS-RA in its setting. These
observations motivate measuring held-out, unshaped success or return under fixed environment-step and
wall-clock budgets, with comparison count and latency reported separately. They do not yield a cost-adjusted
performance bound. % ledger: 8, 9, 12, 19

\newpage
\section{Limitations and open questions}

P's displayed method leaves the inactive-agent reward mask, cross-coalition embedding, pooling rule, and
critic targets unresolved, so the conditional identities cannot diagnose its realised updates. Repeat
comparator calls may be correlated, and the toy variance calculation is a frozen-policy diagnostic. We have
no controlled finite-training experiment, no measured comparison-to-return sensitivity, and no theorem of
improved learned-policy return. % ledger: 10, 11, 12, 15, 16, 19

A direct audit would trace a coalition entry or reentry through P's reward tensor, inactive-agent mask, and
critic target. Matched-critic, neutral-potential, and scale-matched random-potential controls could then
compare held-out unshaped outcomes at fixed step and time budgets while counting calls. Whether a ranking
effect survives, and whether it pays for those calls, remains open. % ledger: 9, 12, 17, 19, 20

\bibliographystyle{plainurl}
\bibliography{references}
\end{document}
